% TOP_PROBLEM: {"problemId": "problem.furstenberg-no-bigeodesics-conjecture-planar-fpp", "canonicalTitle": "Furstenberg No-Bigeodesics (Planar FPP, Standard Minimum Moment)", "releaseRank": 485, "primaryDomain": "probability_ergodic_dynamics", "primaryDomainLabel": "Probability, ergodic theory & dynamics", "displayStatus": "open_with_solved_subcases", "releaseStatus": "open_with_solved_subcases"}
% !TeX program = lualatex
% ATTEMPT_STATUS: unresolved
% Model: GPT-6 Astra Ultra; continuation dated 27 September 2026.
\documentclass[11pt]{article}
\usepackage[margin=25mm]{geometry}
\usepackage{fontspec}
\setmainfont{Latin Modern Roman}
\usepackage{amsmath,amssymb,mathtools}
\usepackage{unicode-math}
\setmathfont{Latin Modern Math}
\usepackage{xurl,hyperref}
\hypersetup{colorlinks=true,urlcolor=blue}
\setcounter{secnumdepth}{0}
\setlength{\parindent}{0pt}
\setlength{\parskip}{5pt}
\setlength{\emergencystretch}{3em}
\begin{document}
\section*{\texorpdfstring{Furstenberg No-Bigeodesics (Planar FPP, Standard Minimum Moment)}{Furstenberg No-Bigeodesics (Planar FPP, Standard Minimum Moment)}}\label{485}
\subsection{Definitions and mathematical statement}
Assign independent edge times with a non-atomic law $G$ on $[0,\infty)$ to the unoriented nearest-neighbor edges of ${\mathbb{Z}}^2$. Assume
\[
 {\mathbb{E}}\bigl[\min(t_1,t_2,t_3,t_4)^2\bigr]<\infty
\]
for four independent samples from $G$. For vertices $x,y$, define $T(x,y)$ as the infimum of the sum of edge times over finite paths from $x$ to $y$. A bigeodesic is a doubly infinite self-avoiding nearest-neighbor path $(v_i)_{i\in{\mathbb{Z}}}$ satisfying
\[
 \sum_{k=i}^{j-1}t_{\{v_k,v_{k+1}\}}=T(v_i,v_j)
 \quad\text{for every }i<j.
\]
The assertion is that the product probability of existence of any bigeodesic is zero. The moment condition concerns the minimum of four times, not an individual edge time. No prescribed directions, density for $G$, curvature of a limit shape, or exponential moment is assumed. This selected version is not equated with a separately stated moment-free version.
\par\smallskip\textit{Formulation and concept references:} \href{https://arxiv.org/pdf/1512.00804}{[S1]}.

\subsection{Short English statement}
Under the specified non-atomic minimum-moment assumption, is every doubly infinite path in planar first-passage percolation almost surely nonoptimal somewhere along its length?
\subsection{Research attempt}

\paragraph{Scope and primary-source audit (27 September 2026).}
This GPT-6 Astra Ultra attempt uses independent, identically distributed
nonnegative edge times with a non-atomic law. Non-atomicity does not mean
that the law has a density. The only moment assumption in the target is
the displayed second moment of the minimum of four samples. In particular,
an individual edge may have infinite expectation.

The supplied Damron--Hanson manuscript [S1], Section 2.1, assumption A1',
uses exactly that minimum moment. Its Theorem 1 also assumes geometric
regularity at selected points of the limit shape and concludes absence of
bigeodesics with an end in the resulting fixed sector. The simpler
Theorem 1.1 in its introduction uses an individual first moment and a
differentiable boundary. Neither form is an unconditional no-bigeodesic
theorem under only this entry's assumptions.

Ahlberg--Hanson--Hoffman [E1], Theorem 1.1 and Section 2.4, assumption A1,
prove that the probability a distant specified vertex belongs to an
infinite geodesic from the origin tends to zero uniformly in that vertex.
Their A1 again permits the present minimum moment. This is an important
input for the discussion below, but it is a statement about a specified
vertex; the rate needed by a naive union over a growing boundary is not
part of its conclusion. No integrable directed last-passage theorem is
substituted for undirected first-passage percolation here.

The main attack is to formulate an exact sequence of finite-box events
whose probability limit is the probability of a bigeodesic through the
origin. This avoids assuming directions for an unknown random path.
The work also proves the basic existence and uniqueness facts directly,
checks the moment distinction on an explicit heavy-tailed law, and rules
out every straight coordinate line by independent local improvements.
The remaining probability estimate is isolated after these reductions.

\paragraph{The minimum moment does not imply an individual first moment.}
Let a nonnegative random variable $t$ have survival function
\[
 \mathbb P(t>s)=(1+s)^{-3/4},\qquad s\geq0.
\]
This is a continuous law with no atom, but
$\mathbb E t=\int_0^\infty(1+s)^{-3/4}\,\mathrm ds=\infty$.
For four independent copies and $M=\min(t_1,t_2,t_3,t_4)$,
\[
 \mathbb P(M>s)=(1+s)^{-3},\qquad
 \mathbb E M^2=2\int_0^\infty\frac{s}{(1+s)^3}\,\mathrm ds=1.
\]
More generally, the exponent $\alpha$ in $(1+s)^{-\alpha}$ gives
\[
 \mathbb E M^2=\frac{2}{(4\alpha-1)(4\alpha-2)}
                        \quad\text{if }\alpha>1/2.
\]
Thus replacing the target's moment by $\mathbb E t<\infty$ would discard
valid examples. The elementary arguments below mostly do not use any
moment assumption at all; this makes their limits particularly transparent.

\paragraph{Geodesic foundations with no hidden moment input.}
Since $G$ has no atom at zero, every edge time is strictly positive almost
surely, simultaneously for the countably many lattice edges. It is finite
because the law is supported on $[0,\infty)$, not on an extended interval
including infinity. Finite loops can therefore be erased at a strict
decrease in passage time.

\textbf{Lemma 485.1 (long inexpensive paths are exponentially unlikely).}
There are constants $c>0$ and $0<\rho<1$ such that, for every vertex $x$
and positive integer $n$,
\[
 \mathbb P\bigl(\exists\text{ a self-avoiding $n$-edge path from $x$
                       with time at most }cn\bigr)
 \leq\frac43\rho^n.
\]

\emph{Proof.}
Choose $\delta>0$ so small that
$0<q=\mathbb P(t\leq\delta)<1/36$. This is possible because the
distribution function is continuous, starts at zero and tends to one.
Set $c=\delta/2$ and $\rho=6\sqrt q<1$.
If a path of $n$ distinct edges has time at most $\delta n/2$, at least
$\lceil n/2\rceil$ of its edges have time at most $\delta$.
For a fixed self-avoiding path, independence and a union over subsets of
its edges bound this probability by
$2^n q^{\lceil n/2\rceil}\leq2^n q^{n/2}$.
There are at most $4\cdot3^{n-1}$ self-avoiding paths of length $n$
starting at $x$: four possibilities for the first step, and at most
three thereafter because immediate reversal is forbidden. A union bound
gives the stated estimate.\hfill$\square$

By summability and Borel--Cantelli, for each $x$ there is an almost surely
finite $N_x$ such that every self-avoiding path from $x$ with at least
$N_x$ edges has time greater than $c$ times its length. Intersect these
probability-one events over all vertices. This controls excessive
excursions of a prospective minimizer without any expectation of $t$.

\textbf{Corollary 485.2.} Almost surely, for every pair of vertices there
exists a unique finite geodesic between them; and from every vertex there
exists at least one singly infinite geodesic.

\emph{Proof.}
For fixed $x,y$, take any deterministic finite path between them and call
its finite time $C$. A self-avoiding competing path of time at most $C$
has fewer than $\max\{N_x,C/c+1\}$ edges by the preceding bound.
There are finitely many such paths. Their minimum is attained, and loop
erasure shows it is the infimum over all finite paths.

For uniqueness, two different simple paths with the same fixed endpoints
have different undirected edge sets. Condition on all edge times except
one edge in their symmetric difference. Equality of the two passage
times then specifies a single value of that remaining edge. Independence
and non-atomicity make its conditional probability zero. There are only
countably many finite paths, so almost surely no such tie occurs for any
pair simultaneously. This applies to finite induced boxes as well.

To obtain an infinite geodesic from $x$, let endpoints tend to infinity
and consider their unique finite geodesics from $x$. Their lengths tend
to infinity. There are finitely many first steps; retain an infinite
subsequence with the same first step. Repeat for the second step, and
so on, taking the diagonal limit. Every finite segment in the resulting
infinite path eventually occurs in a finite geodesic and is minimizing.
The limiting path is self-avoiding because any repeated vertex would
already appear in one of its finite prefixes.\hfill$\square$

The last result produces one-sided paths, not bigeodesics. Taking two
independent choices of one-sided geodesics at the origin does not imply
that their concatenation is minimizing across the origin. That missing
cross-origin optimality is exactly what the following finite events test.

\paragraph{An exact finite-box reformulation of a root bigeodesic.}
Let
\[
 B_n=\{z\in\mathbb Z^2:\|z\|_\infty\leq n\},\qquad
 \partial B_n=\{z:\|z\|_\infty=n\}.
\]
Write $T_n$ for the distance using only edges whose endpoints lie in
$B_n$. Define $F_n$ to be the event that there are distinct
$a,b\in\partial B_n$ and a $T_n$-geodesic from $a$ to $b$ passing
through the origin. This event depends on finitely many edge times.
By positivity and uniqueness, it can be checked by
\begin{equation}\label{ra485:finiteevent}
 F_n=\bigcup_{\{a,b\}\subset\partial B_n}
 \{T_n(a,b)=T_n(a,0)+T_n(0,b)\}.
\end{equation}
If equality holds, concatenating the two shortest paths gives a shortest
walk. It cannot have a positive-cost loop, so it is a path through zero.

\textbf{Proposition 485.3 (nested crossing criterion).}
For positive finite edge times,
\[
 F_{n+1}\subseteq F_n,
 \qquad
 \{\text{there is a global bigeodesic through }0\}
                          =\bigcap_{n\geq1}F_n.
\]
Uniqueness of finite geodesics is not needed for this proposition.

\emph{Proof.}
Suppose a box geodesic witnessing $F_{n+1}$ is rooted at its visit to
zero. Follow each of its two arms from zero up to its first visit to
$\partial B_n$. The concatenated truncated segment lies in $B_n$.
It is shortest among paths in $B_{n+1}$ between its new endpoints,
because a shorter competitor could replace that segment in the original
geodesic. Hence it is also shortest among the smaller collection of
paths contained in $B_n$, and witnesses $F_n$.

A global bigeodesic through zero supplies the same truncated segment in
every box. Each arm eventually reaches the box boundary, since a
self-avoiding infinite path cannot remain in a finite set. A global
minimizer is also a minimizer when competitors are restricted to a box,
so all $F_n$ occur.

Conversely suppose every $F_n$ occurs. Choose one witnessing path in
each box and root it at zero. Each arm has at least $n$ edges because
its endpoint has sup norm $n$. Orient the pair of arms by any fixed
ordering of the two neighbors of zero used by the path. For each $k$
there are only finitely many possible pairs of $k$-edge initial arms.
Repeated subsequence extraction gives a limiting doubly infinite path
$(v_i)_{i\in\mathbb Z}$ with $v_0=0$. Every fixed finite rooted
segment agrees with segments of witnesses in boxes of arbitrarily
large size. In particular the limit is self-avoiding.

Fix $i<j$ and any finite competing path from $v_i$ to $v_j$ in the
whole lattice. That competitor lies in some box $B_m$. For a sufficiently
large selected witness, its segment indexed from $i$ to $j$ agrees with
the limit, and both it and the competitor lie in its box. The witness
is shortest there, so the passage time of the limiting segment is no
greater than the competitor's time. Taking the infimum over all finite
competitors proves the bigeodesic property.\hfill$\square$

Consequently
\begin{equation}\label{ra485:probabilitytarget}
 \mathbb P(\text{a bigeodesic visits }0)
                  =\lim_{n\to\infty}\mathbb P(F_n).
\end{equation}
By translation invariance, vanishing of this one probability implies
absence of every bigeodesic: the existence event is the union over
countably many vertices of the event that one visits that vertex.
Conversely absence certainly forces the probability at zero to vanish.
Thus \eqref{ra485:probabilitytarget} is an exact reduction of the target,
without directions, shape curvature, or moments stronger than those given.

The use of box distances is legitimate only because the exhaustion
argument eventually tests every finite global competitor. For a fixed
$n$, a $T_n$-geodesic need not be a global geodesic: leaving the box may
produce a shortcut. Neither one finite witness nor witnesses whose
centers escape to infinity establish a bigeodesic through a fixed vertex.

\paragraph{Trying to close the reduction with available sparsity.}
Let $\mathcal T_x$ denote the union of all singly infinite global
geodesics starting at $x$. The theorem in [E1] gives
\[
 q_n:=\sup_{\|v\|_\infty=n}\mathbb P(v\in\mathcal T_0)
                              \longrightarrow0.
\]
If a bigeodesic visits zero, one of its arms first reaches some
$v\in\partial B_n$. Starting at $v$ and following the other direction
along that bigeodesic gives an infinite geodesic containing zero. Thus
\[
 \{0\text{ lies on a bigeodesic}\}
       \subseteq\bigcup_{v\in\partial B_n}\{0\in\mathcal T_v\}.
\]
Translation invariance changes the probability in each term to
$\mathbb P(-v\in\mathcal T_0)$. Since $|\partial B_n|=8n$,
\begin{equation}\label{ra485:unionloss}
 \mathbb P(0\text{ lies on a bigeodesic})\leq8nq_n.
\end{equation}
The verified result $q_n\to0$ does not make this upper bound tend to
zero. It would suffice to prove $q_n=o(n^{-1})$, but that stronger rate
is not established here or supplied by the quoted theorem.

Equivalently, the expected number of boundary vertices $v$ whose infinite
geodesic tree contains zero is $o(n)$. A bigeodesic only requires a
bounded number of suitable arms at each scale; sublinear expectation
does not exclude that possibility. Turning this sparsity into absence
requires more geometric information than the direct union estimate.

The finite formulation has a similar obstruction. There are
$\binom{8n}{2}$ possible endpoint pairs in
\eqref{ra485:finiteevent}. A probability tending to zero for each
specified pair is not enough to bound their union. A uniform
$o(n^{-2})$ bound would suffice by the union bound, but is again a
stronger proposed input, not a consequence established in this attempt.
The finite and infinite formulations involve different events, so their
probabilities are not silently substituted for one another.

\paragraph{A complete restricted exclusion by local detours.}
\textbf{Proposition 485.4.} Almost surely no horizontal or vertical
coordinate line is a bigeodesic.

\emph{Proof.}
Non-atomicity ensures that the law is not concentrated at a single value.
Choose finite $0<a<b$ with
$r=\mathbb P(t\leq a)>0$ and $s=\mathbb P(t\geq b)>0$.
Take an integer $m$ so large that $mb>(m+2)a$.
Consider a segment of $m$ horizontal edges on a fixed horizontal line,
and the parallel segment one row above, joined to its endpoints by two
vertical edges. Require all $m$ original edges to have time at least
$b$, and all $m+2$ detour edges to have time at most $a$.
The two edge sets are disjoint, so this event has probability
$s^m r^{m+2}>0$. On it, the original line segment is not minimizing.

Place infinitely many such rectangles with disjoint edge sets along the
line, separated by at least one unused column. Their improvement events
are independent. The probability that none of the first $N$ improvements
occurs is $(1-s^m r^{m+2})^N$, which tends to zero. Hence the fixed line
fails the bigeodesic condition almost surely. Apply the same argument
to vertical lines and take the countable intersection over all coordinate
lines.\hfill$\square$

This proof works even when the support is a very short interval: one
chooses $m$ large enough to offset the two extra detour edges. It does
not require zero weights, large gaps in the support, or a density for
the law. Its limitation is that the tested line was fixed before the
environment was sampled. A random optimizing path preferentially avoids
bad rectangles, so those independent events cannot simply be reused
along a path selected from the same weights.

\paragraph{Direction and quantifier stress tests.}
Even a theorem excluding one prescribed direction at a time need not
exclude a path with a random exceptional direction. The logical issue
already occurs for a uniform random variable $U\in[0,1]$:
$\mathbb P(U=\theta)=0$ for every deterministic $\theta$, while the
random set $\{U\}$ is nonempty with probability one. An uncountable
union of null events need not be null. Nor has this attempt proved that
an arbitrary infinite geodesic possesses a limiting direction. The
finite-box criterion avoids both assumptions but still needs its own
probability estimate.

There is also no inference from unique finite geodesics to the absence
of bigeodesics. On an infinite tree with strictly positive edge lengths,
the unique combinatorial path between two vertices is always minimizing,
and every doubly infinite simple path is a bigeodesic. This example is
outside the planar square-lattice model; it shows precisely that local
uniqueness alone cannot be the missing principle. The ability to make
competing detours in the square lattice must be exploited on the random
path itself.

\paragraph{Finite verification and outcome.}
The companion exact-integer program checks
\eqref{ra485:finiteevent} by shortest-path computations in coupled boxes
of radii one through four. It uses explicitly generated positive rational
weights with distinct sums for distinct simple edge sets, implemented
by an integer part and binary perturbations. These deterministic finite
environments are not asserted to be samples from the infinite non-atomic
product law. Among the 24 tested environments, the numbers with a root
crossing were respectively $24,22,21,19$ at radii $1,2,3,4$.
All 72 successive nestedness comparisons passed, and each distance-additivity
witness was checked to contain zero in its reconstructed shortest path.
The full output is recorded separately; no observed finite frequency is
extrapolated to \eqref{ra485:probabilitytarget}.

The script also verifies the heavy-tail moment formula and the detour
inequality using rational parameters. All stochastic claims used in the
proofs come from the explicit independence, countability, and summability
arguments above, not from those finite computations.

\textbf{Outcome: UNRESOLVED.} The work proves finite geodesic existence
and uniqueness, existence of one-sided geodesics, the exact nested
finite-box criterion, and exclusion of every coordinate line. These are
elementary reductions and benchmarks, not a new no-bigeodesics theorem.
The infinite-geodesic sparsity theorem [E1] is an external established
input whose limitation in the direct union argument is quantified by
\eqref{ra485:unionloss}.

The first missing statement is
$\mathbb P(F_n)\to0$ under the original minimum-moment hypothesis.
Concrete possible next steps are a multiscale argument that bounds this
root-crossing event without paying for every boundary pair, a stronger
control of the random set of geodesic ancestors than its expected size,
or a detour construction that remains valid after the candidate path
has been chosen from the environment. None is established here.
No prescribed-direction, finite-mean, or shape-regularity condition is
added to the unresolved target.

\subsection{Sources}
\begin{itemize}
\item[S1]Furstenberg No-Bigeodesics (Planar FPP, Standard Minimum Moment); exact editorial selection using the primary formulations and hypotheses recorded in the source receipts.
\url{https://arxiv.org/pdf/1512.00804}.
\textit{Formulation source.}
\item[E1] D. Ahlberg, J. Hanson, and C. Hoffman,
\emph{The number of geodesics in planar first-passage percolation grows
sublinearly}, arXiv:2208.11576v1, 24 August 2022,
Theorem 1.1 and Section 2.4, assumption A1.
\url{https://arxiv.org/pdf/2208.11576}.
\textit{Primary author manuscript; exact minimum-moment scope and uniform
specified-vertex conclusion verified 27 September 2026.}
\item[E2] M. Damron and J. Hanson, \emph{Bigeodesics in first-passage
percolation}, arXiv:1512.00804v2, 1 July 2016,
Theorem 1.1, assumption A1', and Theorem 1.
\url{https://arxiv.org/abs/1512.00804v2}.
\textit{Bibliographic identification of supplied source [S1]; distinguishes
the selected minimum moment from the extra geometric hypotheses in the
proved directional theorem. Consulted 27 September 2026.}
\end{itemize}
\end{document}
